\begin{examplebox}
\textbf{例1：两根竖段已知一根，求另一根和面积（直接应用）}

椭圆 $C:\frac{x^2}{16}+\frac{y^2}{4}=1$，
$A_1(-4,0),A_2(4,0)$。按本结论构造 $P_1,P_2$。
若 $|A_1P_1|=1$，求 $|A_2P_2|$ 与梯形面积。

\textbf{解：} $a=4,b=2$，所以
\[
|A_1P_1|\,|A_2P_2|=b^2=4.
\]
故 $|A_2P_2|=4/1=4$，
\[
S=a(d_1+d_2)=4(1+4)=\boxed{20}.
\]
本椭圆的面积最小值是 $2ab=16$；此例 $20>16$，合理。

\medskip
\textbf{例2：逆向构造切线与切点（可迁移题型）}

仍取 $a=4,b=2$，已知 $P_1=(-4,1)$、$P_2=(4,4)$。
问 $P_1P_2$ 是否与椭圆相切？若相切，求切点 $P$。

\textbf{解：} 有符号高度 $y_L=1,y_R=4$，
$y_Ly_R=4=b^2$，满足逆向相切条件。
\[
k=\frac{4-1}{2\cdot4}=\frac38,
\qquad m=\frac{4+1}{2}=\frac52.
\]
故切线是 $y=\frac38x+\frac52$。
利用切点公式 $x_0=-a^2k/m$、$y_0=b^2/m$，
\[
x_0=-\frac{16\cdot(3/8)}{5/2}=-\frac{12}{5},
\qquad y_0=\frac4{5/2}=\frac85.
\]
所以
\[
\boxed{P\left(-\frac{12}{5},\frac85\right)}.
\]
验算：$(-12/5)^2/16+(8/5)^2/4=9/25+16/25=1$。

\medskip
\textbf{例3：给切点，求面积并理解最值（两种捷径）}

椭圆 $\frac{x^2}{16}+\frac{y^2}{4}=1$ 上
$P=(2\sqrt3,1)$，求该切线所围成的梯形面积。

\textbf{解法A（不必求交点）：} 已知 $y_0=1$，
直接使用 $S=2ab^2/|y_0|$，
\[
S=\frac{2\cdot4\cdot4}{1}=\boxed{32}.
\]
\textbf{解法B（核对）：} 切线为
\[
\frac{2\sqrt3}{16}x+\frac14y=1
\ \Longleftrightarrow\ y=4-\frac{\sqrt3}{2}x.
\]
故 $y_L=4+2\sqrt3$，$y_R=4-2\sqrt3$，
\[
y_Ly_R=16-12=4=b^2,
\quad S=4(y_L+y_R)=4\cdot8=\boxed{32}.
\]
对比：在短轴端点 $P=(0,2)$ 处切线为 $y=2$，
所得图形是宽 $8$、高 $2$ 的矩形，面积 $16$，恰好最小。

\medskip
\textbf{例4：易混辨析——把端点切线换成准线会怎样？}

椭圆 $\frac{x^2}{16}+\frac{y^2}{4}=1$，
半焦距 $c=\sqrt{16-4}=2\sqrt3$，
两条\textbf{准线}为
\[
x=\pm\frac{a^2}{c}=\pm\frac8{\sqrt3},
\]
而\textbf{长轴端点处切线}是 $x=\pm4$，两者不同。
仍取例3的切线 $y=4-\frac{\sqrt3}{2}x$。
它与 $x=\pm4$ 相交的有符号高度积是 $4$；
但与两条准线相交的高度分别为
\[
4+\frac{\sqrt3}{2}\cdot\frac8{\sqrt3}=8,
\qquad
4-\frac{\sqrt3}{2}\cdot\frac8{\sqrt3}=0.
\]
两条准线上的相应竖直距离乘积为 $8\cdot0=\boxed0$，
\textbf{并非} $b^2=4$。因此不能把“端点切线”偷换成“准线”。

\medskip
\textbf{例5：高价值变式——为什么只有 $x=\pm a$ 锁住定值？}

椭圆 $\frac{x^2}{a^2}+\frac{y^2}{b^2}=1$ 的切线
$y=kx+m$ 分别交两条对称竖线 $x=-d$、$x=d$
（$d>0$）于 $M,N$。求这两个交点的\textbf{有符号纵坐标积}。

\textbf{解：} 因为切线条件为 $m^2-a^2k^2=b^2$，
\[
y_M=m-dk,\qquad y_N=m+dk,
\]
\[
\boxed{y_My_N=m^2-d^2k^2=b^2+(a^2-d^2)k^2.}
\]
当 $d=a$，斜率项完全抵消，积恒为 $b^2$；
当 $d\ne a$，积一般随 $k$ 变化。
若竖线在椭圆外侧（$d>a$），两个高度还可能异号，
此时\textbf{不能}把有符号乘积当成两段长度乘积。

\medskip
\textbf{例6：只用一个高度做最值比较（迁移练习）}

若切点 $P(x_0,y_0)$ 满足 $|y_0|=b/2$，
不求 $x_0$，直接计算 $S$ 与 $S_{\min}$ 的比值。

\textbf{解：} 利用 $S=2ab^2/|y_0|$，
\[
S=\frac{2ab^2}{b/2}=4ab,
\qquad S_{\min}=2ab,
\]
故 $\boxed{S/S_{\min}=2}$。
\textbf{迁移点：} 面积只由切点纵坐标的绝对值决定；
同一水平线上两个对称切点，得到相同的面积。
\end{examplebox}
